% concatenated from bndgen3.tex
\documentclass[11pt,reqno]{amsart}




\usepackage{amsmath,amssymb,amsthm,enumitem,xcolor,tikz,tikz-cd,subcaption,
  hyperref,mathtools,float,graphicx}
\usepackage{stackengine,scalerel}
% \usepackage[style=alphabetic,sorting=nty,doi=false,isbn=false,url=false]
%{biblatex}
%\renewcommand*{\bibfont}{\small}


\pgfdeclarelayer{fg}
\pgfsetlayers{main, fg}

%\addbibresource{refs.bib}

\usepackage{cleveref}

\numberwithin{equation}{section}
\tolerance=1000
\hypersetup{  pdfborder={0 0 0},
        colorlinks   = true,
        urlcolor     = blue,
        linkcolor    = blue,
        citecolor   = red}

\renewcommand{\ll }{\langle\hspace{-.7mm}\langle }
\newcommand{\rr }{\rangle\hspace{-.7mm}\rangle }
\makeatletter
\newcommand\bigcdot{\mathpalette\bigcdot@{.5}}
\newcommand\bigcdot@[2]{\mathbin{\vcenter{\hbox{\scalebox{#2}{$\m@th#1\bullet$}}}}}
\newcommand{\overbar}[1]{\mkern 0.8mu\overline{\mkern-0.8mu#1\mkern-0.8mu}\mkern 0.8\
mu}
\makeatother
\newcommand{\Cc}{\mathbb{C}}
\newcommand{\N}{\mathbb{N}}
\newcommand{\Z}{\mathbb{Z}}
\newcommand{\R}{\mathbb{R}}
\newcommand{\Q}{\mathbb{Q}}
\newcommand{\E}{\mathbb{E}}
\newcommand{\F}{\mathbb{F}}
\DeclareMathOperator{\St}{St}
\DeclareMathOperator{\Sym}{Sym}
\DeclareMathOperator{\C}{C}
\DeclareMathOperator{\cent}{Z}
\DeclareMathOperator{\Out}{Out}
\DeclareMathOperator{\Aut}{Aut}
\DeclareMathOperator{\Core}{Core}
\DeclareMathOperator{\lcm}{lcm}

\let\Im\relax
\DeclareMathOperator{\Im}{Im}

\newcommand{\f}{ \leqslant_{f}}
\newcommand{\vr}{ \leqslant_{vr}}
\newcommand{\nvr}{ \lhd_{vr}}
%\newcommand{\n}{\unlhds}

\stackMath
\def\unlhds{\ThisStyle{\mathrel{%
      \stackinset{r}{.75pt+.15\LMpt}{t}{.1\LMpt}{\rule{.3pt}{1.1\LMex+.2ex}}
      {\SavedStyle\leqslant}%
}}}


\newcommand{\comment}[1]{\footnote{\textcolor{red}{#1}}}
\newcommand{\blue}[1]{\textcolor{blue}{#1}}

\newcommand{\leqs}{\leqslant}
\newcommand{\geqs}{\geqslant}
\newcommand\restr[2]{{% we make the whole thing an ordinary symbol                   
                \left.\kern-\nulldelimiterspace % automatically resize the bar with \
\right                                                                               
                #1 % the function                                                    
                \littletaller % pretend it's a little taller at normal size          
                \right|_{#2} % this is the delimiter                                 
}}

\newcommand{\littletaller}{\mathchoice{\vphantom{\big|}}{}{}{}}
\newcommand{\ov}{\overline}
\newcommand{\ot}{\widetilde}

% THEOREM Environments ---------------------------------------------------           
\newtheorem{thm}{Theorem}[section]
\newtheorem{prop}[thm]{Proposition}
\newtheorem{lemma}[thm]{Lemma}
\newtheorem{cor}[thm]{Corollary}
\newtheorem{claim}{Claim}[section]

\theoremstyle{definition}
\newtheorem{defn}[thm]{Definition}
\newtheorem{question}[thm]{Question}
\newtheorem{problem}[thm]{Problem}
\newtheorem{constr}[thm]{Construction}

\theoremstyle{remark}
\newtheorem{ex}[thm]{Example}
\newtheorem{rem}[thm]{Remark}
\newtheorem{exercise}[thm]{Exercise}
\newtheorem{notation}[thm]{Notation}


\title{Rank 1 phenomena for 1-relator and deficiency 1 groups}
\author[Button]{J.\,O.\,Button
}
\address{Selwyn College, University of Cambridge,
Cambridge CB3 9DQ, UK}
%\email{jb}

\keywords{bounded, generation\
}
\subjclass[2020]{20E06, 20F65}

\begin{document}
\vspace*{-1cm}

\begin{abstract}
  We examine second bounded cohomology and mod p homology in finite index
  subgroups of 1-relator groups and groups with a presentation of
  deficiency at least one. We use this to
  determine exactly which 1-relator groups are boundedly generated, as well
  as the groups of deficiency at least one
  up to a class of groups that conjecturally
  do not exist.
\end{abstract}


        \maketitle

    \setcounter{tocdepth}{1}
\tableofcontents

%%%%%%%%%%%%%%%%%%%                             %%%%%%%%%%%%%%%%                     
%%%%%%%%%%%%%%%%%%%           New Section       %%%%%%%%%%%%%%%%                     
%%%%%%%%%%%%%%%%%%%                             %%%%%%%%%%%%%%%%                     

\section{Introduction}

It has long been known that in semisimple Lie groups, rank 1 lattices and
lattices of higher rank behave very differently, even just when considering
their group theoretic properties. For instance, if we were to take the
groups $G=SL_2(\Z)$ and $H=SL_n(\Z)$ for any $n\geqs 3$, we could
distinguish $G$ from $H$ in many ways. To give a range of examples, $H$ has
distorted infinite cyclic subgroups unlike $G$, but $G$ is large (has a
finite index subgroup possessing a surjective homomorphism to a  non-abelian
free group) whereas $H$ has property (T) and all of its non-trivial
normal subgroups have finite index. As for the asymptotic behaviour of
the number of subgroups of a given finite index (see \cite{lssg} for
details about this), $G$ has the fastest possible subgroup growth for
a finitely generated group and this is of growth type $n^n$, whereas
$H$ has subgroup growth of type $n^{\log{n}/\log{\log{n}}}$. Subgroup growth
is related to profinite and pro-$p$ properties of a group. Whilst neither
$G$ nor $H$ themselves have interesting pro-$p$ behaviour for any prime $p$,
any finite index subgroup $L$ of $H$ has a small pro-$p$ completion,
meaning that $L_{\hat{p}}$ is $p$-adic analytic (see Section
3 for definitions) whereas there exist finite index subgroups
of $G$ for which this is not true.

Continuing with more properties, whilst it is the case that neither $G$ nor
$H$ have any non-trivial surjective homomorphisms to $\R$ (although many
finite index subgroups of $G$ do), $G$ has an infinite dimensional
vector space of quasimorphisms whereas $H$ does not. Finally $H$ has the
property of bounded generation (see Section 6) whereas $G$ is far away from
that.

Given we have so many properties distinguishing rank 1 lattices from those of
higher rank and that these properties are purely group theoretic, we can
take a family $\mathcal F$ of commonly occurring finitely presented groups
along with one of our distinguishing properties above. We can then try to see if
all groups in $\mathcal F$ have this property.
If so then we can claim that
the groups in $\mathcal F$ have either rank 1 or higher rank behaviour.
However the evidence for this particular property might be inconclusive
in that some groups in $\mathcal F$ could have this property and some
do not. Also inconclusive but more compelling would be if some groups
have the property whereas it is open for other groups but no counterexamples
are known. This would then start to point towards rank 1 or higher rank
behaviour for $\mathcal F$. One could then take another of our properties
and see if the same pattern emerges.

This was undertaken in \cite{flm} for $\mathcal F$ being all of the
the mapping class groups of closed orientable connected surfaces of genus
$g$ where $g\geqs 1$. The point there was that it could be argued mapping
class groups sometimes exhibit rank 1 behaviour and sometimes that of
higher rank. However the focus in that paper was on finding rank 1 properties
satisfied by all groups in $\mathcal F$. Those properties were referred to
as rank 1 phenomena and we will adapt that term here. One of their rank 1
phenomena was that of our group $G$ having a finite index subgroup $L$
such that $L_{\hat{p}}$ is not $p$-adic analytic, with Theorem
1.4 of that paper stating that mapping class groups have this rank 1 phenomenon.
This is then used to obtain Corollary 1.5 which says that the profinite
completion of a mapping class group is not boundedly generated and thus mapping
class groups themselves are not boundedly generated.

However another way to see that mapping class groups are not boundedly
generated is that they all have an infinite dimensional space of
homogeneous quasimorphisms (see Section 2 for details),
for instance because they are acylindrically hyperbolic
groups. In this paper we will look at two particular rank 1 phenomena
for a given finitely presented group $G$: the first is that $G$ has an
infinite dimensional space of homogeneous quasimorphisms, which under
the finite presentation hypothesis is equivalent to having infinite
dimensional second bounded cohomology. Our second rank 1 phenomenon is
that $G$ has some finite index subgroup $H$ where there is a prime $p$
such that the pro-$p$ completion $H_{\hat{p}}$ is not $p$-adic analytic.

This then brings us to ask which families of commonly occurring finitely
presented groups might satisfy one or both of our rank 1 phenomena. Our
focus for this paper is the well studied class of 1-relator groups. We
will also consider the class of groups with a
presentation of positive deficiency, that is where we have more generators
than relators. (Incidentally the papers \cite{ltt} and \cite{hll}
classify all groups $G$ of
positive deficiency which actually are lattices in connected Lie groups:
it turns out that $G$ is a lattice in one of 
$PSL(2,\R)$, $PSL(2,\Cc)$, $\R\times PSL(2,\R)$, $\R$ or $\R\times\R$.)
Note that this class generalises that of 1-relator groups,
at least if we ignore the cyclic groups coming from 1-generator
1-relator presentations. Indeed any naturally occurring class of groups
often contains exceptions like these (eg the family $\mathcal F$ above
did not contain the mapping class group of the genus 0 surface).
Our viewpoint here will be that
the exceptions are groups which are too small to assign rank 1 or higher
rank behaviour to them. In the positive deficiency case these small
groups are ($\Z$ and) the soluble Baumslag - Solitar groups $BS(1,n)$.
So now we can ask whether all non-small groups of positive 
deficiency satisfy either of our rank 1 phenomena. It is known
(see Sections 2 and 3) that both
of these hold for any group with deficiency at least two (here there are
no small groups). However for groups of deficiency one, it is unknown for
second bounded cohomology and false for the pro-$p$ phenomenon. But
we can then ask is it true that at least one of these rank 1 phenomena
holds for all of our groups?

In this paper we show that this is indeed true for the family of
non-small 1-relator groups. We also show that this holds for all
non-small groups of positive deficiency, with one potential class
of exceptions which are conjectured not to exist. The outline of
the paper is that we look at when a group has infinite dimensional
second bounded cohomology in Section 2. We are quickly able to use known
results to obtain Corollary \ref{toz} which deals with all groups $G$ of
positive deficiency except when $G$ is an ascending HNN extension of some
finitely generated group $H$.
In Section 3 we introduce our other rank 1 phenomenon concerning pro-$p$
completions of finite index subgroups of $G$. Again deficiency at least two
is easily taken care of and for deficiency one groups, this phenomenon
takes on a particularly tractable form in Definition \ref{vsa} which we
name virtually sufficient homology.

When $G$ is an HNN extension of $H$,
we divide into two cases: where $G=H\rtimes\Z$ and when $G$ is a strictly
ascending HNN extension of $H$. The first case is easily dealt with in
Section 4. We have from known results that $H$ must be free and here
a range of different arguments are available to show that $G$ has both
of our rank 1 phenomena. The heart of the paper is in Section 5 on
strictly ascending HNN extensions. The first problem here is that we
do not now know $H$ will be free, hence the potential
exceptions in the deficiency one case although \cite{bi} Section 2
conjectures that this does not occur. But even if $H=F_r$ is free
of finite rank $r$, the family $\mathcal S$ of
strictly ascending HNN extensions of $F_r$
can be harder to deal with than the free by cyclic groups $F_r\rtimes\Z$.
Although some general results are known for groups $G$ in
$\mathcal S$, for instance
\cite{bs} and \cite{mut}, a particularly intractable case is when 
$G$ contains a Baumslag - Solitar group $BS(1,n)$ for $|n|\geq 2$ whereupon
the geometry of $G$ will not be well behaved. To give a particular
example, if $F_2$ is free on $a,b$ and we consider the injective endomorphism
$\theta$ of $F_2$ sending $a$ to $a^2$ and $b$ to $b$ then there are clearly
elements growing both polynomially and exponentially under $\theta$ but
the strictly ascending HNN extension $\langle t,a,b\rangle$ of $F_2$ will
not be hyperbolic relative to any collection of proper subgroups. However
any free-by-cyclic group $F_r\rtimes_\alpha\Z$ having elements growing both
polynomially and exponentially under the automorphism $\alpha$ will be
relatively hyperbolic.

Therefore we proceed by showing in Theorem \ref{main} that for $r\geq 2$,
any ascending HNN extension of $F_r$ containing $BS(1,n)$ for $|n|\geqs 2$
possesses our second rank 1 phenomenon of virtually sufficient homology.
In doing this, we adapt an argument in \cite{lgv3} to show in Theorem
\ref{melg} that if an injective endomorphism of $F_r$ sends some
non-identity element $w\in F_r$ to a conjugate of $w^r$ then $w$ is a
power of a primitive element.

So we conclude that one or other of our rank 1 phenomena holds for every
non-small group of positive deficiency one (away from the exceptional
case that is not known to occur). However in Section 6 we can apply
this to the property of bounded generation, because either of our rank 1
phenomena holding for $G$ implies that $G$ is not boundedly generated.
Our final section is on 1-relator groups. Here this exceptional case
cannot occur, so one or other of our two rank 1 phenomena does hold for
every non-small 1-relator group. Hence Corollary 7.2 which states that the
boundedly generated 1-relator groups are precisely the cyclic groups and
the soluble Baumslag - Solitar groups $BS(1,n)$. We finish by investigating
what happens in the 1-relator case
if we strengthen the second bounded cohomology condition to
being acylindrically hyperbolic. Now we would expect that any non-cyclic
1-relator group which is neither acylindrically hyperbolic nor has virtually
sufficient homology would be isomorphic to $BS(m,n)$ for $m,n$ coprime.
But we find all the exceptional 1-relator groups $G$ in Theorem \ref{last}
and show it can happen that $G$ is a generalised Baumslag - Solitar group
but not a Baumslag - Solitar group.

The author would like to thank F.\,Fournier - Facio for help with references.


\section{Bounded cohomology}

Here we are interested in the second bounded cohomology and its connection
to quasimorphisms. We give brief details and
definitions because this is an area that is very well covered in the
literature: see the papers cited below and references therein for more detail.

Given any group $G$, let $H^2_b(G;\R)$ be the second bounded cohomology of
$G$, which is an $\R$-vector space. We have the comparison map $c$ which
is a linear map from $H^2_b(G;\R)$ to the usual cohomology $H^2(G;\R)$.
If $Q_h(G)$ is the vector space of homogeneous quasimorphisms
$q:G\rightarrow\R$ (meaning that there exists $D\geqs 0$ such that
$|q(gh)-q(g)-q(h)|\leqs D$ and $q(g^n)=nq(g)$ for all $g,h\in G$ and
$n\in\Z$) and $H^1(G;\R)$ is the vector space of genuine homomorphisms
then the kernel $EH^2_b(G;\R)$ of $c$ (where $E$ stands for exact) can
be identified with the quotient vector space $Q_h(G)/H^1(G;\R)$.
Thus if $G$ is finitely generated and the vector space $Q_h(G)$ has infinite
dimension then so does $H^2_b(G;\R)$. As for the converse,
it is conceivable that there exists a finitely
generated group with a finite dimensional
space of homogeneous quasimorphisms but with $H^2_b(G;\R)$ infinite
dimensional. However if $G$ is finitely presented the the usual cohomology
$H^2(G;\R)$ has finite dimension and so here the converse does hold.


The free group $F_n$ for $n\geq 2$ was initially shown to have infinite
dimensional second bounded cohomology.
As for other groups with this
property, if $G$ has a surjective homomorphism to a group $H$ with infinite
dimensional second bounded cohomology then so does $G$.
It is also the case that if $G$ has infinite dimensional second
bounded cohomology then so do the finite index subgroups of $G$.
However it is not known whether the finite index supergroups of $G$ do.


More generally
by \cite{epfuj}, having infinite dimensional second bounded cohomology
holds for any non-elementary word hyperbolic group. This also holds for
relatively hyperbolic groups (which here will mean any group hyperbolic
relative to some collection of proper subgroups) and even
acylindrically hyperbolic groups, which is essentially in \cite{bstfgn}
despite preceding the definition of acylindrical hyperbolicity. These
results are obtained by analysing a suitable action of our group $G$
on some hyperbolic space. This suggests that further progress can be
made by specialising to actions on simplicial trees. This is indeed
the case as can be seen in \cite{sstocan}, \cite{hffmstrs}. We will just
use two earlier results here, one for HNN extensions and one for amalgamated
free products.
\begin{thm} (\cite{kj2} Theorem 1.2), see also \cite{grig})
\label{kj2hnn}
  Suppose that $G$ is any group that splits as an HNN extension
  (which is equivalent to $G$ having a surjection to $\Z$). If $G$ is
  an HNN extension of $A$ with edge groups $C, \phi(C)$ both
  properly contained in $A$ then $G$
  has infinite dimensional second bounded cohomology.
\end{thm}

  (Note that any $G$ satisfying these hypotheses must contain $F_2$.)
  \begin{cor} \label{toz}
    If $G$ is any finitely presented group with a surjection to $\Z$ but
    with finite dimensional second bounded cohomology
    then $G$ is an ascending HNN extension
    of a finitely generated group.
  \end{cor}
  \begin{proof} Any HNN extension $G$ which is finitely presented can be
    ``folded'', which involves changing the HNN extension but not $G$ itself,
    to an HNN extension of some group $A$ with edge subgroups
    $C$ and $\phi(C)$ of $A$
    where $A$ and $C$ (thus $\phi(C)$ too) are finitely generated
    (see any of \cite{bistr} or \cite{bsh} or \cite{dnwdybk} VI Theorem 4.5
    for instance). If neither $C$
    nor $\phi(C)$ are equal to $A$ then $G$ has infinite dimensional
    second bounded cohomology by Theorem \ref{kj2hnn}. Otherwise we are now in
    the conclusion of the corollary.
   \end{proof}       
\begin{thm} (\cite{kj2} Theorem 1.1), see also \cite{grig})
\label{kj2amg}
Suppose that $G$ is any group that splits as an amalgamated free
product $A*_CB$.  If $|C\backslash A/C|\geq 3$ and $|B/C|\geq 2$
(these indices are allowed to be infinite) then $G$
  has infinite dimensional second bounded cohomology.
\end{thm}
(Again any $G$ satisfying these hypotheses must contain $F_2$.)

Two useful corollaries follow from this.
\begin{cor} (\cite{kj2} Corollary 1.2), see also \cite{grig})
\label{kj2inf}
Suppose that $G$ is any group that splits as an amalgamated free
product $A*_CB$.  If $A$ is infinite but $C$ is finite with
$|B/C|\geq 2$ then $G$
  has infinite dimensional second bounded cohomology.
\end{cor}
\begin{cor} (\cite{kj2} Corollary 1.3), see also \cite{grig})
\label{kj2abl}
Suppose that $G$ is any group that splits as an amalgamated free
product $A*_CB$.  If $A$ is abelian, $|A/C|\geq 3$ and $|B/C|\geq 2$
then $G$ has infinite dimensional second bounded cohomology.
\end{cor}

Note: care is sometimes needed, given that we
cannot necessarily pass up to finite
index supergroups. In particular if a finitely generated group 
is virtually free (and not virtually cyclic)
then, as pointed out in the proof of \cite{kj2} Theorem
1.3, we see it has infinite dimensional second bounded cohomology because
it is (non-elementary) hyperbolic, not because it has $F_n$ as a finite
index subgroup. For another example, take any group $G$
with a finite presentation
of deficiency at least two (so at least two more generators than
relators) By the famous result
\cite{bp} of Baumslag and Pride, $G$ is large which means it has a finite
index subgroup possessing a surjective homomorphism to $F_2$. However
this does not directly tell us that $G$ has infinite dimensional second
bounded cohomology so we would need to use
Corollary \ref{toz} for groups of deficiency at least two;
indeed by adapting this argument
\cite{mofup} tells us that these groups are 
not just virtually acylindrically hyperbolic but actually acylindrically
hyperbolic (this is another example of a property which is not known
to pass to finite index supergroups).

For examples of groups which have finite dimensional second bounded cohomology,
this is true for all amenable groups. Our viewpoint here will be that such
groups are too ``small'' to be able to assign rank 1 or higher rank
phenomena to them. The other main source are nearly all
$S$-arithmetic subgroups of linear algebraic groups, in particular
groups such as $SL(n,\Z)$ for $n\geq 3$ (but not $n=2$ as that group is
virtually free).


We will require one application of the above results here. A {\bf generalised
Baumslag - Solitar group} is the fundamental group of a finite graph
of groups where each edge and vertex group is a copy of $\Z$. We also
have {\bf generalised Baumslag - Solitar groups of rank $n$}, where we
  replace $\Z$ with $\Z^n$ for some fixed $n$ in the definition. Although
  we are only use the rank 1 case here, we will assume for the next result
  that the definition includes the rank $n$ case, as the proofs are the same
  for arbitrary $n$.

  \begin{thm} \label{gbs}
    Let $G$ be a generalised Baumslag - Solitar group with underlying graph
    $\Gamma$. Assume without loss
    of generality that the finite graph of groups is reduced; that is the
    inclusion of any edge group into any adjoining vertex group is proper
    unless this edge is a self loop.\\
    (i) $G$ is never acylindrically hyperbolic.\\
    (ii) $G$ does have infinite dimensional second bounded cohomology
    unless $\Gamma$ is a single point, or $\Gamma$ is a single
    edge with index 2 inclusions on both sides, or $\Gamma$ is a single
    loop with (at least) one inclusion of the edge group being equal to the
    vertex group.
 \end{thm}   
 \begin{proof}
   The without loss of generality comment in the statement holds because we can
   perform an elementary collapse, that is we can contract any edge which is
   not a self loop if its edge group is equal to a vertex group on one side.
   We then relabel the edge inclusions on the other side, so that we still
   have a generalised Baumslag - Solitar group, and this does not change
   the fundamental group.

   To show $G$ is not acylindrically hyperbolic, note that any vertex group
   $V$ (which is a copy of $\Z^n$) is commensurable with $gVg^{-1}$ for an
   arbitrary $g\in G$, because the inclusions are all finite index, so the
   associated Bass - Serre tree has finite valency but $V$ and $gVg^{-1}$
   are both vertex stabilisers. In the language of \cite{motr} $V$ is
   an $s$-normal subgroup of $G$ but Theorem 3.7 (c) there tells us that any
   $s$-normal subgroup of an acylindrically hyperbolic group has to be
   acylindrically hyperbolic itself.

   For (ii), note that whenever we form an HNN extension or a non-trivial
   amalgamation $A*_CB$ (namely when the edge group $C$ is a proper subgroup
   of both $A$ and $B$) then the resulting group contains the vertex group(s)
   with infinite index. We now consider the generalised Baumslag - Solitar
   group $G$ case by case.

   If the defining graph $\Gamma$ is not a tree then to form our fundamental
   group $G$ we take a maximal tree in $\Gamma$ and contract each edge one
   by one, forming amalgamations each time, until we are left with a single
   point with self loops whereupon we create an HNN extension with each
   loop. Thus unless $\Gamma$ has only one vertex, when we take the first
   self loop in order to form the first HNN extension, we have that at both
   ends of this loop, the edge group is contained in the vertex group with
   infinite index. This property will then continue to hold as we move through
   the other self loops. Thus we are done by Theorem \ref{kj2hnn} unless
   we had one vertex and we are also done if $\Gamma$ consisted of one
   vertex with more than one self loop.

   If $\Gamma$ is just a single self loop then we are again done by
   Theorem \ref{kj2hnn} unless one of the inclusions is equal to the
   vertex group. In other words $G$ is just an ascending HNN extension
   of the soluble group $\Z^n$, so will itself be soluble and thus amenable.

   Otherwise $\Gamma$ is a tree. In this case we can intersect all of
   the finitely many vertex and edge subgroups of $\Gamma$ 
   to obtain a subgroup $N$ contained in each vertex group $V_i$ and
   each edge group $E_{ij}$ with finite index.
   As these are all abelian, $N$ is normal in $V_i$
   and in $E_{ij}$. Thus $N$ is
   normal in $G$ itself because the vertex groups generate $G$,
   since here $\Gamma$ is a tree since no HNN extensions are formed.
   Thus $G/N$ can be seen as the fundamental group of a finite graph of
   groups with
   vertex groups $V_i/N$ and edge groups $E_{ij}/N$. But all vertex and
   edge groups in $G/N$ are finite, hence $G/N$ is
   virtually free. Therefore $G/N$, thus $G$,
   has infinite dimensional second bounded
   cohomology by the comments above, unless $G/N$ is virtually
cyclic and so $\Gamma$ is a point (whereupon
   $G=\Z^n$), or a single edge with index 2 inclusions on either side
   (in which case $G=\Z^n*_C\Z^n$ is soluble) or $\Gamma$ is a single self
   loop with one inclusion equal to the vertex group $\Z^n$, so that $G$
   is an ascending HNN extension of $\Z^n$ and so is again soluble.
   \end{proof}

   In the case $n=1$, our three exceptions for $G$ will be $\Z$, the Klein
   bottle group and the soluble Baumslag - Solitar groups $BS(1,n)$ (which
   gives the Klein bottle again when $n=-1$). It is well known that these
   are precisely the virtually soluble groups of deficiency (at least) one.
    
    
   
\section{Profinite and pro-$p$ techniques}

   From the last section, we see that the mapping
   class group $\mathcal{M}$ (say of a closed orientable surface with genus
   $g\geqs 1$) has infinite dimensional
   second bounded cohomology, because it 
   is acylindrically hyperbolic due to its action on the
   curve complex. This can be viewed as a rank 1 phenomenon, in that it
   is shared with rank 1 lattices of semisimple Lie groups but not
   with higher rank lattices. Another rank 1 phenomenon for $\mathcal{M}$
   is given in \cite{flm} Theorem 1.4 and can be thought of as a feature
   of its finite index subgroups. This is that $\mathcal{M}$ has a finite
   index subgroup $H$ whose pro-$p$ completion $H_{\hat{p}}$ is not a
   $p$-adic analytic group. The background to this is
   explained in \cite{lssg} which is a wide ranging summary of results in this
   area. In particular, for any prime $p$ the $p$-adic analytic groups are
   exactly the pro-$p$ groups of finite Pr\"ufer
   rank, whereupon their subgroup
   growth is polynomial in the index $n$. This allows one to formulate
   an analog of the congruence subgroup property for abstract finitely
   generated groups which again can be regarded as a rank 1 phenomenon,
   as mentioned at the start of
   \cite{lssg} Chapter 7. The idea is that small subgroup
   growth, meaning of growth type strictly less than $n^{\log n}$, is a
   higher rank
   phenomenon, whereas subgroup growth of type strictly bigger than $n^{\log n}$
     is a rank 1 phenomenon. The connection to $p$-adic analytic groups
     is through the famous Golod - Shafarevich inequality.

     As explained in
     \cite{lssg} 4.6, suppose we have a finite presentation
     $\langle X\,|\,R\rangle$ for an abstract group $G$, so here the
     deficiency of this presentation is $|X|-|R|$. For any prime $p$,
     this can be regarded as a pro-$p$ presentation for the pro-$p$
     completion $G_{\hat{p}}$. If $d_p$ is the minimum number of generators
     for $G_{\hat{p}}$ then there exists a pro-$p$ presentation
     $\langle Y_p\,|\,S_p\rangle$ for $G_{\hat{p}}$ where $|Y_p|=d_p$ and the
     deficiency $|Y_p|-|S_p|$ is the same as our original deficiency
     $|X|-|R|$. Now the Golod - Shafarevich inequality is that
     $d_p^2/4\leq |S_p|$ for every pro-$p$ presentation
     $\langle Y_p\,|\,S_p\rangle$ where $|Y_p|=d_p$. Then \cite{lssg}
     Theorem 4.6.4 states that if $G_{\hat{p}}(\ncong\Z_p)$ has an open
     subgroup violating this inequality then $G_{\hat{p}}$ has subgroup
     growth which is strictly bigger than $n^{\log{n}}$, so $G_{\hat{p}}$ is
     not $p$-adic analytic because its subgroup growth is bigger than
     polynomial. But then the subgroup growth of $G$ is also 
     strictly bigger than $n^{\log n}$ so that we have our rank 1
     phenomenon for $G$.

     However it might happen, as in the mapping class group
     (when the genus is at least 3) that a group $G$ has a very small
     pro-$p$ completion for all primes $p$ (say if $G$ is perfect or
     the abelianisation $G/G'$ is cyclic) but $G$ has a finite index subgroup
     $H$ where $H_{\hat{p}}$ is not pro-$p$ analytic for some prime $p$. In
     that case the larger subgroup growth type for $H$ also holds for $G$.
     
     

     For a finitely generated group $G$ we have that $d_p$ is equal to the
     rank of $H_1(G,\F_p)$ (say by \cite{gz2} Proposition 4.1.3). For
     groups with positive deficiency, this allows us to detect this
     rank 1 phenomenon in a very direct way.  (This is based
     on \cite{lbsgcg} which deals with lattices in $SL(2,\Cc)$.
     See also \cite{lac} Theorem 5.1 for a similar result on
     finitely generated groups $G$ in terms of the rank of both
     $H_1(G,\F_p)$ and $H_2(G,\F_p)$.)
     \begin{thm} \label{padan}
       Suppose  we have a group $G$
       which is defined by a presentation $\langle X\,|\,R\rangle$
       of deficiency $|X|-|R|$ at least 1. Then the following are
       equivalent:\\
              (i) There exists a finite index subgroup $H$ of $G$ with
       the abelianisation $H/H'$ needing at least 3 generators.\\
       (ii) there exists a finite
       index subgroup $L$ of $G$ and a prime $p$ such that the pro-$p$
       completion $L_{\hat{p}}$ is not $p$-adic analytic\\
       (iii) The supremum of rank($H_1(S,\F_p)$) over all finite index subgroups
       $S$ of $G$ is infinite.\\
       If $|X|-|R|\geq 2$ then these conditions always hold.
\end{thm}
       \begin{proof}
         Suppose (i) holds then $H$ will also have a presentation of
         deficiency at least 1
     (say by the Reidemeister - Schreier process),
         and there will be some $p$
         such that the rank of $H_1(H,\F_p)\geq 3$.
         Then on applying the above comments to $H$, we obtain a pro-$p$
         presentation $\langle Y_p\,|\,S_p\rangle$ for $H_{\hat{p}}$ where
         $|Y_p|-|S_p|$ is equal to the deficiency of $H$ which is at least
         1 and $d_p=H_1(H,\F_p)\geq 3$.
         Thus we violate the Golod - Shafarevich result above, so
         $H_{\hat{p}}$ is not $p$-adic analytic for this prime $p$.

      If (ii) holds then $L_{\hat{p}}$ has infinite Pr\"ufer rank. This means
      that the number of topological generators needed for an open subgroup
      $P$ of $L_{\hat{p}}$ is unbounded as we vary over all such $P$.
      But each $P$ corresponds to an abstract
      finite index subgroup $S$ of $L$ which is subnormal in $L$ and with
      the rank of $H_1(S,\F_p)$ equal to that of $H_1(P,\F_p)$, thus our
supremum is infinite.

Then it is clear that (iii) implies (i). In the case of deficiency
at least 2, the Golod - Shafarevich inequality never holds so (ii) is
satisfied by $G$ itself.
\end{proof}

     Thus we can think of any group $G$ in Theorem \ref{padan}
     as satisfying a different rank 1 phenomenon to that of having
     infinite dimensional bounded second cohomology from the last section.
     Note that because of Theorem \ref{padan}, it makes sense to allow for
     finite index subgroups of $G$ rather than expecting $G$ itself to
     satisfy these conditions. For instance any group $G$ defined by a
     2-generator, 1-relator presentation will have $G/G'$ needing at most
     2 generators.

     We now formalise our rank 1 phenomenon in the case of deficiency 1
     presentations by using the following definition.

     \begin{defn} \label{vsa}
       A group $G$ defined by a deficiency 1 presentation
       has {\bf virtually sufficient abelianisation} if there is a finite index
       subgroup $H$ with $H/H'$ having rank at least 3.
     \end{defn}

     Certainly not all deficiency 1 groups will have virtually sufficient
     abelianisation, e.g. $\Z$ and $\Z\times\Z$. Indeed $BS(m,n)$ has
     virtually sufficient homology if and only if $m$ and $n$ are coprime,
     for instance see \cite{me1r} Section 4
     where many other examples of 2-generator 1-relator
     groups without virtually sufficient homology are given. (It might be
     noted here that all those other exceptions seem to be acylindrically
     hyperbolic; we return to this point in the last section.)
          However one case which
     does satisfy this definition is when our group $G$ is linear in
     characteristic zero but not virtually soluble. Then (for any $p$)
     the supremum of $H_1(H,\F_p)$ over all finite index subgroups $H$
     of $G$ is infinite (see \cite{lssg} Window 9, Corollary 18). Thus any
     deficiency 1 group which is not virtually soluble but which is linear
     in characteristic zero has virtually sufficient homology.
     Finally
     we have mentioned that this rank 1 phenomenon implies subgroup growth
     that is strictly bigger than $n^{\log{n}}$. But in the deficiency
     1 case, \cite{lac} Theorem 1.5 (ii) states that if $G$ has virtually
     sufficient abelianisation then its subgroup growth will be a
     lot faster, in fact not far off exponential growth $n^{n/\log{n}}$.
     

     
     \section{Free by cyclic groups}

     We have seen that the only candidates for deficiency 1 groups $G$ which
     do not have infinite dimensional second bounded cohomology are
     those which are an ascending HNN extension of a finitely generated
     group. This can be divided into two cases, where the ascending HNN
     extension is strict which we deal with in the next section, and
     otherwise whereupon $G$ is of the form $K\rtimes\Z$ for $K$
     the finitely generated kernel of our map to $\Z$ in Corollary \ref{toz}.
     (If the first Betti number of $G$ is at least 2 then it is possible for
     $G$ to be in both cases, but we would then apply the stronger results
     in this section to $G$.)
     
Here we define a free by cyclic group as a group of the form
$F_r\rtimes\Z$ where $F_r$ is a free group of finite rank $r$. We
allow any $r$, including $\Z$ for $r=0$, as well as
$\Z\times\Z$ and the Klein bottle
group for $r=1$. These last three exceptional groups do not satisfy either of
our two rank 1 phenomena but this is due to them being ``small'' groups.
All free by cyclic groups have a deficiency 1 presentation. Conversely
a result in \cite{dss} tells us that if $G$ is of the form
$K\rtimes\Z$ for $K$ finitely generated and $G$ has a deficiency 1
presentation then $K$ is free. If $r\geq 2$ then there are many ways to
see that any such $G$ has both of our rank 1 properties. However
one method we
cannot use for virtually sufficient abelianisation is by linearity, as
it is unknown if all free by cyclic groups are linear. Here we proceed
by using mainly acylindrical hyperbolicity for the bounded cohomology
and a short argument for the virtually sufficient abelianisation.

\begin{thm} \label{fbyz}
  Any free by cyclic group $G=F_r\rtimes \Z$ with $r\geq 2$
  has infinite dimensional second bounded cohomology and virtually
  sufficient abelianisation.
\end{thm}  
\begin{proof} For the first part, we are done if $G$ is hyperbolic,
  relatively hyperbolic or acylindrically hyperbolic. By \cite{ghst}
  or \cite{dm?}, we will find that $G$ is acylindrically hyperbolic
  unless the defining automorphism $\alpha$ of $F_r$ has finite order
  in $Out(F_r)$ whereupon $G$ is virtually $F_r\times \Z$. Although virtually
  is not useful to us here, we note that $G$ is definitely not acylindrically
  hyperbolic, as if it were then $F_r\times\Z$ would be too.
Instead we argue that by \cite{lvrf} Proposition 4.1 we see $G$ is
a generalised Baumslag - Solitar group (here we mean the $r=1$ case).
Thus by Theorem \ref{gbs} and the fact that $G$ is not soluble, we
conclude that $G$ has infinite dimensional second bounded cohomology.

For virtually sufficient abelianisation, we use the trick
(as in \cite{lac} Proposition 8.7 or \cite{mearx} Lemma 3.6)
that $\alpha$ gives rise to the abelianised automorphism
$\overline{\alpha}$ of $\Z^r$ and hence to the mod $p$ abelianised
automorphism $\overline{\alpha}_p$ of $\F_p^r$. As the mod $p$ automorphism
has finite order $d$ say, we can take the $d$-fold cyclic cover $H$ of $G$,
which is $\langle t^d,F_r\rangle$ where $t$ is the stable letter
inducing $\alpha$ by conjugation. Then we have that
$(\overline{\alpha}_p)^d=\overline{(\alpha^d)}_p$ is the identity. But $H$
has index $d$ in $G$ and the rank of $H_1(H,\F_p)$ is $r+1$, with the 1
coming from the $\Z$ in the abelianisation of $H$ given by the stable
letter $t^d$. As $r+1\geq 3$, we are done for any prime $p$.
\end{proof}

  

\section{Strictly ascending HNN extensions}

We have now considered both of our rank 1 phenomena for all presentations
of deficiency at least 1, except for when the presentation has deficiency
exactly 1 and the resulting group is an HNN extension $\langle t,A\rangle$
where $t$ is the stable letter, the vertex group $A$ is finitely
generated and the HNN extension is formed from an endomorphism
$\theta:A \rightarrow A$ which is injective but not surjective. The obvious
case where this case results in a deficiency 1 presentation is
when $A$ is a free group $F_r$ of rank at least 1. For $r=1$ we have
the Baumslag - Solitar groups $BS(1,n)$
where $|n|\geqs 2$, which we have seen are soluble and therefore too ``small''
to satisfy either of our rank 1 phenomena.
These particular
Baumslag - Solitar groups are also {\bf non - Euclidean} which we define
to mean Baumslag - Solitar groups $BS(m,n)$ where $|m|\neq |n|$.

Of course if a group is hyperbolic then it cannot contain any
Baumslag - Solitar
group $BS(m,n)$, where $m,n\neq 0$ and we can say without loss of generality
given the symmetries in the presentation that $1\leqs m$ and $m\leqs |n|$.
In fact no ascending HNN extension of $F_r$ can contain $BS(m,n)$
if $m\geqs 2$, in which case $BS(m,n)$ contains $\Z\times\Z$ anyway, 
so the obstruction to hyperbolicity here is the existence of subgroups
isomorphic to $BS(1,n)$.
In \cite{mut} the converse was shown for strictly ascending HNN extensions.
Note that for these groups $G$ the presence of a
Baumslag - Solitar subgroup is equivalent to the existence of a {\bf periodic
  conjugacy class} for the endomorphism $\theta:F_r\rightarrow F_r$;
namely we have elements $x\neq e$ and $g$ in $F_r$ and
$d\in\Z\setminus\{0\},i>0$ such that $\theta^i(x)=gx^dg^{-1}$
(so that it is the conjugacy class of $x$ in $F_r$ which is periodic).
Thus we do have infinite dimensional second bounded cohomology for $G$ if 
there are no periodic conjugacy classes for $\theta$.

However strictly ascending HNN extensions of free groups $F_r$ for
$r\geqs 2$ which contain non - Euclidean Baumslag - Solitar
groups seem quite mysterious.
They are all known to be residually finite by \cite{bs}.
However we have the following examples.
\begin{ex} \label{drsp}
  Following \cite{drusap}, given $r\geqs 1$ and $d\neq 0$ we define
  the {\bf Dru\c{t}u - Sapir group} $DS(r,d)$ to be the group
  with presentation
  \[\langle t,a_1,\ldots ,a_r\,|\,ta_1t^{-1}=a_1^d, \ldots ,ta_rt^{-1}=a_r^d
    \rangle\]
  which is clearly an ascending HNN extension of $F_r$ and is strictly
  ascending if and only if $|d|\geqs 2$. If $|d|=1$ or $r=1$ then it is
  a linear group but it is shown in \cite{drusap} that otherwise it is
  non linear. Moreover consider the ascending HNN extension $C(r,d)$
  of $F_r$ with presentation
  \[\langle t,a_1,\ldots ,a_r\,|\,ta_1t^{-1}=a_2, \ldots ,
    ta_{r-1}t^{-1}=a_r,ta_rt^{-1}=a_1^d
    \rangle.\]
  Note that $C(r,d)$ has a 2-generator 1-relator presentation, given by
  eliminating $a_r$ then $a_{r-1}$ and so on, right up to $a_2$. But the
  obvious cyclic cover of $C(r,d)$ having index $r$ is $DS(r,d)$, so
  that $C(r,d)$ is also not linear if $r\geqs 2$ and $|d|\geqs 2$.
\end{ex}

However other ascending HNN extensions of $F_r$ for $r\geqs 2$
can contain non - Euclidean Baumslag - Solitar groups but still be
linear, for instance in \cite{caldun} where in $SL(2,\Cc)$ they take
$t$ and $F_6=\langle a\rangle *F_5$, with $t$ conjugating $F_6$ into itself
such that $tat^{-1}=a^3$.
The Dru\c{t}u - Sapir groups are clearly not hyperbolic and are not
relatively hyperbolic, though they are acylindrically hyperbolic if
$r\geqs 2$ and $|d|\geq 2$.


We now turn to asking which strictly ascending HNN extensions $G$ of finitely
generated free groups have virtually sufficient homology. Note that the
argument used in the last section for free by cyclic groups will not work in
general now. Indeed if the defining endomorphism $\theta$ sends every
generator of $F_r$ into the commutator subgroup $F_r'$ then the induced
mod $p$ abelianised endomorphism of $\F_p^r$ will be zero for every prime $p$.
On the other hand, if $G$ is word
hyperbolic and the defining endomorphism $\theta$ is irreducible then
\cite{hgws} shows that $G$ is cubulated, hence is large and so has
virtually sufficient homology. However this ``hence'' involves going from
cubulation, namely acting freely and cocompactly on a CAT(0) cube complex,
to being virtually special, namely virtually embedding in a RAAG, and this
requires the full Agol - Wise machinery. We do not know of a more elementary
proof of virtually sufficient abelianisation in this case.
Furthermore if $G$ is hyperbolic but $\theta$ is reducible then $G$ is not
known to be cubulated. Here we do not know of an argument that shows
$G$ has virtually sufficient homology, although hyperbolicity means $G$
satisfies our other rank 1 phenomenon.

By \cite{mut} mentioned above, if $G$ is not hyperbolic then
we can now assume that $G$ contains $BS(1,d)$
for $|d|\geq 1$. Relative or acylindric hyperbolicity will not help us
with the homology of finite index subgroups (and nor will linearity,
given Example \ref{drsp}) so we now divide into
two cases according to whether $G$ contains $BS(1,1)$
(or $BS(1,-1)$ but this itself contains $BS(1,1)$),
or just $BS(1,d)$
for $|d|\geqs 2$. In the former case it was shown in \cite{med1} that
$G$ is large and hence has virtually sufficient homology, because for
any $n$ there will exist a finite index subgroup $H$ of $G$ which
possesses a surjective homomorphism to $\Z^n$. We now deal with the latter
case, although first we review some consequences of the 
Reidemeister - Schreier rewriting process which is itself described
in \cite{ls} Proposition II.4.1.

Suppose we have a group $G$
given by the presentation $\langle X\,|\,R\rangle$
(not assumed finite here) and we want to get a presentation for a subgroup
$H$ (not assumed to be finite index)
of $G$. This can be done by taking
a Schreier transversal $S$ for $H$ in $G$ which
is a right transversal for $H$ in $G$, given as words in $X$,
having the property that any initial subword of $s\in S$
is also in $S$. Note that (unless $S=\{e\}$ in which case $H=G$) some element
of $X$ or its inverse must appear in $S$. For any $g\in G$, let $\overline{g}$ 
be the representative of $g$ in $S$, that is the unique $s\in S$ such that
$Hs=Hg$. We then have that
\[Y:=\{sx(\overline{sx})^{-1}\,|\,s\in S,x\in X\}\]
is (throwing away occurrences of the identity) a generating set
for $H$. As for obtaining a set of relators for $H$,
we can do this by considering for $s\in S$ and $r\in R$ the relator
$srs^{-1}$ which can then be
written in terms of the elements of $Y$ to obtain a relator
in $Y$. Doing this over all $s\in S$ and $r\in R$ gives us a defining
presentation for $H$ with these relators and our generating set $Y$.

We first consider what happens when $G$ is a free group. Here
we say that $g\in G$ is {\bf primitive} in $G$ if it is an element of some
free basis for $G$.

\begin{lemma} \label{prim}
  Let $G$ be a free group and $g$ be a primitive element of $G$. Suppose
  we have a subgroup $H$ of $G$ and $k>0$ such that $g^k\in H$ but
  $g,g^2,\ldots ,g^{k-1}$ is not in $H$. Then $g^k$ is a primitive element
  of $H$.
\end{lemma}
\begin{proof}
Let $X$ be a free basis of $G$ in which $g$ appears. Then for
$0\leqs i<j<k$ we have $Hg^i\neq Hg^j$ so that $e,g,\ldots, g^{k-1}$ is a
partial Schreier transversal. On extending to a full Schreier transversal $S$,
we will find elements of our generating set for $H$ by taking the various
expressions $sx(\overline{sx})^{-1}$ which are not equal to the identity.
When we set $x=g$, we can also take $s=e,g,\ldots ,g^{k-2},g^{k-1}$ in turn.
Whereas $e,\ldots ,g^{k-2}$ all return $e$, for $s=g^{k-1}$ we find that
$g^k\neq e$ is an element of our generating set. To see that this
generating set $Y$ is a free
generating set, note that the original set of relators $R$ for $G$ is empty.
Thus the new set for $H$ using the generating set $Y$ is too.
\end{proof}

So now suppose we have $G$ a strictly ascending HNN extension of the free
group $F_n$ (where $n\geqs 2$) which contains some subgroup $BS(1,d)$ for
$|d|\geqs 2$. We said above that this gives rise to a
periodic conjugacy class, so that there is $i>0$ and
  $w,v\in F_n$ where $w$ is not the identity with $\theta^i(w)=vw^dv^{-1}$.
  Note that if $w=u^k$ is a power of some other element $u\in F$ then
  $(\theta^i(u))^k=(vu^dv^{-1})^k$ and $k$th roots are unique in a free
  group, so we can replace $w$ with $u$ and hence
  assume that $w$ is not a proper power.
  We can also tidy this up by first replacing $\theta$ with $\phi=\theta^i$
  which just changes $G$ to its finite cyclic cover of index $i$
  (which is still an ascending HNN extension).
  We can then replace $\phi$ with $\iota_{v^{-1}}\phi$ for $\iota_{v^{-1}}$
  the inner automorphism which is conjugation by $v^{-1}$ and which does
  not change the group. Renaming
  $\iota_{v^{-1}}\phi$ by $\theta$ and the cyclic cover of $G$ by
  $G$ also, we now have $\theta(w)=w^d$.

  We now show that when $|d|\geq 2$ (though not necessarily when $d=\pm 1$),
  this element $w$ is primitive in $F_n$. This result is in a paper \cite{lgv4}
  of Logan which uses the machinery of Mutanguha in \cite{mut}. Here we
  adapt the more basic idea from Lemma 3.4 in an earlier version \cite{lgv3}
  of Logan's paper.
  \begin{thm} \label{melg}
    Let $\theta: F_n\rightarrow F_n$ be an injective endomorphism of the free
    group $F_n$ for $n\geqs 2$ and let $w\in F_n\setminus\{e\}$ be such
    that $\theta(w)=w^d$ for $|d|\geqs 2$ but $w$ is not a proper power.
    Then $w$ is primitive in $F_n$.
  \end{thm}
  \begin{proof}
    An old and classic theorem of G.\,Baumslag in \cite{bmrn} is as follows.
    Say $g_1,\ldots ,g_n$ are elements of the free group $F$ on
    $x_1\ldots ,x_n$ for $n\geqs 2$
    and $w(x_1,\ldots ,x_n)$ is a word which is not a proper power and not
    primitive when regarded as an element of $F$.
    If we have $g_1,\ldots ,g_n,g\in F$ and
    $k\geq 2$ such that the relation $w(g_1,\ldots ,g_n)=g^k$ holds in $F$
    then the subgroup $S:=\langle g_1,\ldots ,g_n,g\rangle$, which must
    clearly have rank at most $n$, actually has rank less than $n$.

    Let us set $g_i=\theta(x_i)$ for $1\leqs i\leqs n$.
    We have a word $w\in F_n\setminus\{e\}$ with $\theta(w)=w^d$ so we can
    put $g=w(x_1,\ldots ,x_n)^{\pm 1}$ (according to whether $d$ is positive
    or negative) and $k=|d|$ to obtain
    \[\theta(w(x_1,\ldots ,x_n))=w(\theta(x_1),\ldots ,\theta(x_n))
      =w(g_1,\ldots ,g_n)=g^k.\]
    As we said that $w$ is not a proper power, we have that $w$ is primitive
    unless
    $\langle \theta(x_1),\ldots ,\theta(x_n),w\rangle$ has rank less than $n$.
    Now the subgroup $\langle \theta(x_1),\ldots ,\theta(x_n)\rangle$ of $F$ has
    rank $n$ because $\theta$ is injective. However it can happen that 
    throwing in $w$ can reduce the rank. But notice that if $n=2$
    then we would be done, because  $\langle \theta(x_1),\theta(x_2)\rangle$
    is non-abelian, so cannot be made cyclic by adding an extra generator.

    Therefore we invoke a rank reduction argument. The key here is that
    if we have the strictly ascending HNN extension $\langle t,F\rangle$
    where $F$ is finite rank free, defined by the endomorphism $\theta$,
    and $S$ is any finitely generated subgroup of $F$ (not necessarily
    of finite index) of rank $m$ say
    with $\theta(F)\leqs S\leqs F$ then
    $\langle t,F\rangle$ is also naturally the ascending HNN extension
    $\langle t,S \rangle$, because $\theta(S)\leqs \theta(F)\leqs S$, so
    we just take a free basis $y_1,\ldots ,y_m$ for $S$ and express each
    $\theta(y_i)$ in terms of this basis.

    Here we will set $S:=\langle \theta(F),w\rangle$ which is clearly
    finitely generated and contains $\theta(F)$. We saw above that if
    the rank of $S$ is equal to that of $F$ then we are done. But
    otherwise the rank is smaller, so we can replace $F$ with $S$ in
    the above and continue since $S$ also contains $w$,
    with the process terminating at some point
    since we cannot go below rank 2.

    However we are changing the vertex group of the extension each time.
    Suppose our rank does reduce when going from $F$ to $S$ but not at the
    next step. Then our application of Baumslag's result is telling us that
    $w$ is primitive in $S$, not necessarily $F$. But supposing that the
    former case holds (and it will hold after finitely many reductions),
        we are done if we can show that $\theta(w)$ is
    primitive in $\theta(F)$, as we then use injectivity of $\theta$.
    Now $\theta(w)=w^d\in\theta(F)$ and as $k=|d|$,
    we will be done by Lemma \ref{prim} on setting
    $G=S$, $g=w$ which is primitive in $G$
    and $H=\theta(F)$ provided $w,\ldots ,w^{k-1}\notin \theta(F)$. So suppose
    we have $u\in F$ such that $\theta(u)=w^j$ where $1\leqs j<k$. Then
    $\theta(u^k)=w^{jk}=\theta(w^{\pm j})$. By injectivity this would say that
    $u^k=w^{\pm j}$ with $u\notin \langle w\rangle$ so that $w$ is a proper
    power, a contradiction.
    \end{proof}
    
    \begin{thm} \label{main}
      Let $\theta:F_r\rightarrow F_r$ be an injective endomorphism
  of the free group $F_r$ where $r\geqs 2$ and let
  $G=\langle t,F_r\rangle$ be
  the corresponding HNN extension. If $BS(1,d)\leqs G$ for $|d|\geq 2$
  then there is
  a finite index subgroup $H$ of $G$ and a prime $p$ such that the rank
  $H_1(H,\F_p)$ is at least 3.
\end{thm}
\begin{proof}
  We have said that we can replace $G$ with a finite cyclic cover
and we can assume
  that we have $w\in F_r\setminus\{e\}$ with $\theta(w)=w^d$.
By Theorem \ref{melg}, $w$ is
actually a primitive element of our free group.
  Thus on taking a free basis $t,a,a_2,\ldots ,a_r$ for $F_r$ with $w=a$ as its
  first element, we obtain for $G$ the presentation
\[G=\langle t,a,a_2, \ldots ,a_r\,|\,tat^{-1}=a^d,
  ta_2t^{-1}=w_2,\ldots ,ta_rt^{-1}=
  w_r\rangle\]
where $|d|\geqs 2$ and $w_2,\ldots ,w_r$ are elements of $F_r$ that will not
particularly concern us.
The abelianisation of $G$ maps onto $\Z$ by sending
$t$ to 1 and all other generators to zero. Moreover by abelianising
the first relation, we obtain $(d-1)\overline{a}=0$ and this vanishes
when we work modulo any prime $p$ dividing $d-1$. If $d=2$ then there are
no such $p$ but we can take the double cyclic cover and on replacing
$t^2$ with $t$, our first relation is now $tat^{-1}=a^4$, hence we can
assume in the above that $d\neq 2$ so such a prime $p$ exists.

Consequently
we have a deficiency 1 presentation for $G$ and a prime $p$ so that, when we
abelianise and reduce modulo $p$, our first relation now vanishes. Hence
we now have $n+1$ variables and $n-1$ equations in $\F_p$ so
$H_1(G,\F_p)$ has rank at least 2. Now given any
finite index subgroup $H$ (say having index $i$)
of $G$, we can obtain a
presentation of $H$ by the Reidemeister - Schreier rewriting process
described above. In this case
$X=\{t,a,a_2,\ldots ,a_r\}$ is our generating set for $G$
and $|R|=r$, thus resulting in a
presentation for $H$ with $ir+1$ generators (throwing away the $i-1$
occurrences of the identity in this process) and $ir$ relators.

This works for any finite index subgroup $H$ of $G$.
The idea now is that if we can find a particular finite index
subgroup $H$ with index $i>1$ and such that
$t,a\in H$ (so in particular $t,a$ will not generate
$G$) then
$x=t$ and $x=a$ are part of our generating set $X$ for $G$. On taking
the identity
$e$ for our element $s$ in our Schreier transversal, we obtain
elements $sx(\overline{sx})^{-1}$ as two generators for $H$, which will be
$t$ and $a$. But there must exist some other generator from
$a_2,\ldots ,a_r$, say $b$, with $b^{\pm 1}$ in $S$ 
because $S$ must contain a length
one element (or $S=\{e\}$ but that is when $H=G$). On now taking
$x=t$ and $x=a$ again but $s=b^{\pm 1}$, we find that
$sx(\overline{sx})^{-1}$ is equal to $t':=b^{\pm 1}tb^{\mp 1}$ and
$a':=b^{\pm 1} ab^{\mp 1}$.
Having obtained all of the $ir+1$ generators for $H$ including these four, we
form the $ir$ relators for $H$ by rewriting the original relators and their
conjugates by elements of $S$ in terms of our new generating set $Y$
for $H$. But our first relator $tat^{-1}a^{-d}$ stays as it is because
$t,a\in Y$. Since $b^{\pm 1}\in S$,
we next conjugate this relator by $b^{\pm 1}$ to obtain
$b^{\pm 1}tat^{-1}a^{-d}b^{\mp 1}$
which will be rewritten as $t'a't'^{-1}a'^{-d}$. It does not matter about
the other relators, because when we now abelianise our deficiency 1
presentation for $H$, we have two relators that disappear modulo $p$
for any prime $p$ dividing $d-1$. Consequently we have that the rank
of $H_1(H,\F_p)$ is at least 3.

We now have to ensure that such an $H$ exists. It certainly will not
if $t,a$ generate $G$. But as we have the relation $tat^{-1}=a^d$
between them, if they did then $G$ would be a quotient of $BS(1,d)$
and so would be metabelian, which is a contradiction as $G$ contains
$F_r$ for $r\geq 2$. So we can hope that there is a finite quotient $G/N$
of $G$ which is not generated by the images $tN$ and $aN$ in $G/N$. Then
$\langle tN,aN\rangle$ will be a proper subgroup of $G/N$ and by
correspondence this will take the form $H/N$, where
$H$ will be a proper finite index subgroup containing $t$ and $a$ so
that we would be done. To obtain such a finite quotient, recall by
\cite{bs} that $G$ is residually finite. Also $G$ is not metabelian, so
that we can take $\gamma \in G''\setminus \{e\}$ and a finite quotient
$Q$ of $G$ in
which $\gamma$ survives. Thus $Q''\neq \{e\}$, but if $Q$ were generated
by the images of $t$ and $a$ then it would be a quotient of
$\langle t,a\rangle$ which is metabelian, so $Q$ would be too and thus
$Q''=\{e\}$.
\end{proof}

We mentioned at the end of Section 3 that \cite{lac} Theorem 1.5 (ii) tells
us that a group with a deficiency 1 presentation having virtually
sufficient homology must have a lower bound for
subgroup growth which is not far off exponential
and in that paper the argument we used in Section 4 is invoked to obtain
this conclusion for all (non-abelian free) by cyclic groups. However if a
finitely generated group is large then it has the highest possible subgroup
growth which is of type $n^n$. Certainly there are deficiency 1 groups which
are not residually finite and which have very small (eg linear) subgroup
growth, as well as the polynomial subgroup growth of the residually
finite groups $BS(1,n)$. As for strictly ascending HNN extensions of finite
rank non-abelian free groups $G$, we have type $n^n$ (via largeness)
for hyperbolic groups with $\theta$ irreducible as well as when $G$
contains $BS(1,\pm 1)$. Theorem \ref{main} above now tells us that we
have at least
near to exponential growth when $G$ contains $BS(1,d)$. Oddly it
is just hyperbolic groups with $\theta$ reducible amongst all strictly
ascending HNN extensions of finite rank non-abelian free groups $G$
where we do not seem to have a sensible lower bound for subgroup growth.

\section{Bounded Generation}

We now mention some applications of our results from the previous sections.
The first is to bounded generation.

\begin{defn} \label{bgdf}
  A group $G$ is {\bf boundedly generated}
  if there exists a finite (generating) set
  $c_1,\ldots ,c_n\in G$ and $N\in\N$ such that
  every $g\in G$ can be expressed as
  \[g=c_{i_1}^{m_1}\ldots c_{i_N}^{m_N}\mbox{ where }
    1\leqs i_1,\ldots ,i_N\leqs n\mbox{ and } m_1,\ldots ,m_N\in\Z.\]
\end{defn}
Another frequently used definition of bounded generation is that $G$
is equal to the product $C_1\ldots C_M$ where
$C_1,\ldots ,C_M$ are cyclic subgroups of $G$ (possibly with repeats).
This clearly implies Definition \ref{bgdf}. For the other way round, note
that if $G$ is boundedly generated by $c_1,\ldots ,c_n$ with word length
$N$ then we can
express $G$ as the product of cyclic groups as long as this product
contains every length $N$ word in $\{1,\ldots ,n\}$. For instance,
if $G$ is boundedly generated by $a,b$ with $N=3$ in Definition \ref{bgdf}
then for
$A=\langle a\rangle$, $B=\langle b\rangle$ we would have
$G=AAABABBBAA$, so that in general we should expect $M$ to be much larger
than $N$.

Obvious examples of finitely generated groups which are boundedly generated
comprise finite groups (but no other torsion groups as there will only be
finitely many possibilities for $g$ in Definition \ref{bgdf}), alongside
abelian groups.
It is also the case that if $G$ is boundedly generated then so are the
finite index subgroups and supergroups of $G$. Moreover bounded generation
is preserved by taking quotients and indeed extensions (meaning that
if $N\unlhds G$ with $N$ and $G/N$ both boundedly generated then so is $G$).
So for instance this means that all virtually polycyclic groups are
boundedly generated.

Slightly less obvious examples can be formed in the
following way. If $G=C_1\ldots C_M$ is boundedly generated and
$\theta:G\rightarrow G$ is an injective (but not necessarily surjective)
endomorphism then on forming the ascending HNN extension
\[\Gamma=\langle G,t\,|\,tgt^{-1}=\theta(g)\mbox{ for }g\in G\rangle,\]
we can move $t$ to the right of any $g$ and $t^{-1}$ to the left,
so that any element in $\Gamma$ can be written as $t^{-p}gt^q$
for $g\in G$ and $p,q\geqs 0$. Hence if $T=\langle t\rangle$ then
$\Gamma=TC_1\ldots C_MT$ is also boundedly generated. In particular
$BS(1,n)$ is boundedly generated. A much less obvious example
is given in the paper \cite{mur} of Muranov. This gives an example of a
finitely generated (though not finitely presented) group which contains
$F_2$ and is a direct limit of a sequence of hyperbolic groups with
respect to a family of surjective homomorphisms.

The other main source of boundedly generated groups are nearly all
$S$-arithmetic subgroups of linear algebraic groups, in particular
$SL(n,\Z)$ is boundedly generated for $n\geq 3$ (but not $n=2$) using
the elementary matrices. To put it loosely, away from small groups
we have that bounded generation is a higher rank phenomenon. Thus
in looking now
for groups which are not boundedly generated, it should be the case that
either of our two rank 1 phenomena are enough to establish this. Indeed
this is well known, as we now summarise:
\begin{prop} \label{ntbg}
  Let $G$ be a finitely generated group.\\
  (i) If $G$ has an infinite dimensional
  space $Q_h(G)$ of homogeneous quasimorphisms then $G$ is not
  boundedly generated. If $G$ is finitely presented and has infinite
  dimensional second bounded cohomology then $G$ is not boundedly
  generated.\\
  (ii) If $G$ is a finitely generated group having a finite index subgroup
  $H$ whose pro-$p$ completion $H_{\hat{p}}$ for some prime $p$ is not
  $p$-adic analytic then $G$ is not boundedly generated.
\end{prop}
\begin{proof}
  The point in (i) is that if $G$ is boundedly generated by $g_1,\ldots ,g_n$
then any homogeneous quasimorphism which is zero on all of
$g_1,\ldots ,g_n$ must be bounded on $G$ and hence
identically zero, by homogeneity. But if we have
$n+1$ linearly independent homogeneous quasimorphisms then we can solve
$n$ linear equations to find a non-trivial homogeneous quasimorphism which
is zero at $g_1,\ldots ,g_n$. Then
$G$ being finitely presented implies that $H^2(G;\R)$ is finite dimensional,
so $H^2_b(G;\R)$ being infinite dimensional is equivalent to $Q_h(G)$
being infinite dimensional.

For (ii), if $G$ were boundedly generated then $H$ would be too, but
the pro-$p$ completion $H_{\hat{p}}$ of a boundedly generated abstract group
$H$ must also be boundedly generated (where we now mean topological
generation). But a boundedly generated pro-$p$ group must be $p$-adic
analytic, say by \cite{lssg} Corollary 12.6.2. Indeed it is further shown
in \cite{pyb} that a boundedly generated abstract group has subgroup
growth of type bounded above by $n^{\log{n}}$.
\end{proof}



We can now apply all of this to groups with a deficiency 1 presentation.
\begin{thm} \label{nbgd1}
  If $G$ is a group defined by a presentation of deficiency 1 then either
  $G$ is isomorphic to $BS(1,n)$ for $n\in\Z$ or
  $G$ is not boundedly generated or $G$ is a strictly ascending HNN extension
  of the finitely generated group $H$ where $H$ is not $FP_2$. Conjecturally
  this last case does not occur.
\end{thm}
\begin{proof}
  There will exist a homomorphism from $G$ onto $\Z$
  which means we can express $G$ as an HNN extension $A*_\phi$.
  We saw in Corollary
  \ref{toz} that we can assume the vertex group $A$ and edge group
  $C$ are finitely generated. If neither $C$ nor $\phi(C)$ are equal
  to $A$ then that corollary tells us we are done by second bounded
  cohomology. If $C=\phi(C)=A$ then $G=A\rtimes\Z$ and we mentioned in
  Section 4 that $K$ will be free, so we are done by Theorem \ref{fbyz}
  using virtually sufficient abelianisation (or by second bounded
  cohomology but virtually sufficient abelianisation is much more
  elementary here). The only exceptions are when $G=\Z$ (which we can
  interpret as $BS(1,0)$), $\Z\times\Z$ or the Klein bottle group
  which are all boundedly generated.

  
  Thus we now know that $C$ is equal to $A$ but $\phi(C)\neq A$ (or
  vice versa but then just swap $\phi(C)$ and $C$).
  so that $G$ is a strictly ascending HNN extension of the finitely
  generated group $A$. If $A$ is of type $FP_2$ then it is known by
  \cite{bi} that $A$ must be free and it is conjectured at the
  end of Section 2 in this paper that this is always the case.

  Hence we are left with $G$ being a strictly ascending HNN extension
  of a finite rank free group. If the rank is
  one then $G=BS(1,n)$ which
  is boundedly generated. If the rank is at least two then
  we are done by second bounded cohomology
  if $G$ is hyperbolic. Otherwise by \cite{mut} $G$ must contain a
  periodic conjugacy class and if this is $\Z\times\Z$ then we are
  done by \cite{med1} using largeness (which can be thought of as
  either using second bounded cohomology on dropping down to a finite
  index subgroup or using virtually sufficient abelianisation).

  So the only case left is when $G$ has a periodic conjugacy class
  giving rise to $BS(1,n)$ for $|n|\geqs 2$, whereupon we are done by 
  Theorem \ref{main} using virtually sufficient abelianisation.
\end{proof}

Looking back at Proposition \ref{ntbg}, we note here that neither
converse is true. For this, we first define the notion of polynomial
index growth. If $G$ is a finitely generated (abstract, profinite
or pro-$p$) group then we can define the index growth function $f_G(n)$
by the supremum of $|Q/Q^{(n)}|$ where $Q$ ranges over all finite quotients
of $G$ and $Q^{(n)}$ is the normal subgroup of $Q$ generated by the $n$th
powers of elements of $Q$. By Zelmanov's resolution of the restricted
Burnside conjecture, $f_G(n)$ is finite and we say that $G$ has
{\bf polynomial index growth} if there exists $C\geq 0$ such that
$f_G(n)\leqs n^C$. Now having bounded generation (in either the abstract,
profinite or pro-$p$ sense) restricts the size of $|Q/Q^{(n)}|$ to being
polynomial in $n$ and so $f_G(n)$ is polynomial in $n$.

For a counterexample to the converse of
(a), we can take any finitely generated residually
finite group $G$ that is virtually soluble but which does not have
finite Pr\"ufer rank (and finitely presented examples exist too).
Then $G$ is amenable thus
has no bounded cohomology but \cite{pysg} shows that $G$ does not have
polynomial index growth, so is not boundedly generated. For (b), examples
of finitely generated groups $G$ (linear in characteristic zero) where
$G$ is not boundedly generated but the profinite completion $\hat{G}$
is were given in \cite{invt}. Thus $\hat{G}$ has polynomial index growth
and so does $G$, because the finite quotients of $\hat{G}$ are the same
as those for $G$. But for any prime $p$ the finite quotients of
$G_{\hat{p}}$ are a subset of this, so $G_{\hat{p}}$ also has polynomial
index growth. Now in the world of pro-$p$ groups, having finite Pr\"ufer
rank, being boundedly generated and having polynomial index growth are
the same thing (see \cite{lssg} Chapter 12 and references within). Thus
$G_{\hat{p}}$ is $p$-adic analytic. We can now repeat with $H$ any finite
index subgroup of $H$, so that $\hat{H}$ will be a finite index subgroup
of $\hat{G}$ and thus will also be boundedly  generated.

We can adapt the heavy machinery used here to give an alternative
argument for Theorem \ref{nbgd1} in the case of strictly ascending HNN
extensions $G$ of finitely generated free groups. It was shown in \cite{brsp2}
as a followup to \cite{bs} that $G$ is virtually residually (finite $p$)
for all but finitely many primes $p$. Now suppose that $G$ is boundedly
generated so that $G_{\hat{p}}$ is also boundedly generated and hence
  is $p$-adic analytic.
  Again by Chapter 12 of \cite{lssg} and related references, this implies that
  $G$ is linear in characteristic zero which in turn implies that
  the rank of $H_1(H,\F_p)$ is unbounded over the finite index subgroups
  $H$ of $G$, so we are done by Theorem \ref{padan}. Note though that
  this does not establish virtually sufficient abelianisation for $G$
  as was achieved in Theorem \ref{main}.
  Rather it says that if $G$ is boundedly
  generated then $G$ both does and does not have virtually sufficient
  abelianisation.
  
  \section{1-relator groups}

  We finish by specialising to 1-relator groups where we would hope that the
  results can be strengthened. We are able to do this for Theorem
  \ref{nbgd1}.

\begin{thm} \label{rel1}
  If $G$ is a finitely generated group defined by a
  1-relator presentation then either
  $G$ is cyclic, is
  isomorphic to $BS(1,n)$ for $n\in\Z\setminus\{0\}$ or
  $G$ has infinite dimensional second bounded cohomology
  or $G$ has a finite index subgroup $H$ where
  $H_{\hat{p}}$ is not $p$-adic analytic for some prime $p$.
\end{thm}
\begin{proof}
  If the presentation for $G$ has three or more generators then we are
  back in the deficiency at least two case, where we know both of
  the latter options are true. If there is one generator then $G$ is
  cyclic. Otherwise we have a 2-generator 1-relator
  presentation and so we can follow the proof of Corollary \ref{toz}.
  However when we write $G$ as an HNN extension with vertex
  group $A$ and edge groups $C,\phi(C)$ which are all finitely generated,
  the Freiheitssatz for 1-relator groups tells us that $C$ and $\phi(C)$
  are free, so that if the HNN extension is ascending then $A$ is free.
\end{proof}

This gives by \ref{ntbg} a result that we have not seen in the literature.
\begin{cor} The boundedly generated 1-relator groups are precisely
  the cyclic groups and the soluble Baumslag-Solitar groups
  $BS(1,n)$.
\end{cor}  

Having seen that all 1-relator groups apart from the ``small'' cases
have at least one of our two rank 1 phenomena, we might want to see
if this still holds on strengthening the first phenomenon from having
infinite dimensional second bounded cohomology to being acylindrically
hyperbolic. A little thought reveals that Baumslag - Solitar groups
of the form $BS(m,n)$ where $m$ and $n$ are coprime will also be
excluded as we know they are not acylindrically hyperbolic and they do
not have virtually sufficient homology. However surprisingly there
are other exceptions.

For this we can use:
\begin{prop} \label{1rel2gen} (\cite{motr} Proposition 4.20) 
Let $G$ be a group with two generators and one defining relator. Then at 
least one of the following holds:
\begin{itemize}
  \item[(i)] $G$ is acylindrically hyperbolic;
  \item[(ii)] $G$ contains an infinite cyclic $s$-normal subgroup.
  More precisely, either $G$ is infinite cyclic or it is an HNN-extension 
of the form $$G=\langle a,b, t\mid  a^t=b, w=1\rangle $$ of a 
2-generator 1-relator group
  $H= \langle a,b \mid w(a,b)\rangle $ with non-trivial center, so that 
$a^r=b^s$ in $H$ for some $r,s \in \Z \setminus\{0\}$. In the latter case $H$  
is (finitely generated free)-by-cyclic and contains a finite index normal 
subgroup splitting as a direct product  
of a finitely generated free group with an infinite cyclic group.
  \item[(iii)] $G$ is isomorphic to an ascending HNN extension
of a finite rank free group.
\end{itemize}
Moreover, the possibilities (i) and (ii) are mutually exclusive.
\end{prop}

In (ii) Theorem C of \cite{krp} tells us that $G$ is a generalised
Baumslag - Solitar group. Thus we need to determine when generalised
Baumslag - Solitar groups do not have virtually sufficient homology. They
certainly do if the group is large. By Theorem 6.7 of \cite{lvrf},
if a generalised Baumslag - Solitar group $G$ is not large (nor cyclic)
then it may be defined by a reduced
graph that is a circle, with edges $e_i$ that
are consistently oriented and with non-zero integer labels $l_i$ and
$r_i$ at the left and the right of $e_i$ and such that $L=\Pi\,l_i$ and
$R=\Pi\, r_i$ are coprime with no $l_i$ or $r_i$ equal to $0$.
Thus if our group is not a Baumslag - Solitar group
then our circle will have 
$n\geq 2$ edges with $l_i,r_i\neq 0,\pm 1$ as the graph is reduced,
whereupon
$G$ will have the deficiency 1 presentation
\[\langle a_1,a_2,\ldots ,a_n,t\,|\,
  a_1^{l_1}=a_2^{r_1}, a_2^{l_2}=a_3^{r_2}, \ldots, a_n^{l_n}
  =ta_1^{r_n}t^{-1}\rangle\]
with $L=l_1\ldots l_n$ and $R=r_1\ldots r_n$ being coprime.


\begin{thm} \label{last} Suppose that
  $G$ is a finitely generated group defined by a
  1-relator presentation with at least 2 generators. If $G$ is neither
  acylindrically hyperbolic 
  nor has a finite index subgroup $H$ where for some prime $p$
  $H_{\hat{p}}$ is not $p$-adic analytic then $G$ is a generalised
  Baumslag - Solitar group with a presentation of the form above.
  Conversely if $G$ has such a presentation then it is neither
  acylindrically hyperbolic nor has a finite index subgroup
  $H$ with $H_{\hat{p}}$ being $p$-adic analytic for some prime $p$.
  \end{thm}
\begin{proof}
  We have now excluded the cyclic case and 3 or more generators follows
  as before. Case (i) in the above gives us acylindrical hyperbolicity
  and Case (iii) is dealt with using the results in Sections 4 and 5
  except when the free group has rank 0 or 1, giving us $BS(1,n)$. Thus
  we are left in Case (ii) whereupon we now know that $G$ is
  a generalised Baumslag - Solitar group which is not large, because of
  its lack of virtually sufficient homology, so $G$ has some presentation
  of the form above.

  We now need to show that all groups defined by such a presentation
(which is of deficiency one)
  fail to have virtually sufficient homology. Note that
  \[a_i^{r_{i-1}\ldots r_1}=a_1^{l_{i-1}\ldots l_1}\mbox{ for any }
    2\leqs i\leqs n\] and
  \[ a_i^{l_i\ldots l_{n-1}l_n}=a_n^{r_i \ldots r_{n-1}l_n}=ta_1^{r_1\ldots r_n}t^{-1}
      \mbox{ for any } 1\leqs i\leqs n.\]
    As $r_{i-1}\ldots r_1$ and $l_i\ldots l_n$ are coprime, we have that
    $a_i\in\langle t,a_1\rangle$ so that $G$ is generated by $t$ and $a_1$.
    Moreover we have shown that the relation $ta_1^Rt^{-1}=a_1^L$ holds
    in $G$ so that $G$ is a quotient of $BS(R,L)$ for $R,L$ coprime. As the
    latter does not have virtually sufficient homology, nor does the
    former. Also $G$ is not acylindrically hyperbolic by
    Theorem \ref{gbs} (i).
\end{proof}
To show there exist groups in Case (ii) which are not genuine
Baumslag - Solitar groups, take $n=2$ in the presentation above so that
we have
\[G:=\langle a,b,t\,|\,
  a^i=b^j, b^k
  =ta^lt^{-1}\rangle\]
where $ik$ is coprime to $jl$ and with none of $i,j,k,l$ equal to $\pm 1$.
Using a ``non-abelian'' Euclidean
algorithm on the coprime integers $j$ and $k$, we can perform Nielsen
moves on our presentation by reducing the powers
$b^j$ and $b^k$ until we are left with $b=u(a,t)$ and $b^m=v(a,t)$ for
some $m$, whereupon we can substitute $u$ into $b^m$, thus
eliminating $b$ and giving us a 2-generator 1-relator presentation
on $a$ and $t$.
By \cite{frdfrm} $G$ is not isomorphic to a Baumslag - Solitar group because
the graph defining $G$ is strongly slide free.

For instance taking
\[G:=\langle a,b,t\,|\,
  a^2=b^3, b^2
  =ta^3t^{-1}\rangle\]
gives us that $b=b^3b^{-2}=a^2ta^{-3}t^{-1}$ so that
$G=\langle t,a\,|\,a^2=(a^2ta^{-3}t^{-1})^3\rangle$.
On setting $u=tat^{-1}$, we see this is an HNN extension
of $H=\langle a,b,u\,|\,a^2=b^3, b^2=u^3\rangle$, namely
$\langle u,a\,|\,a^2=(a^2u^{-3})^3\rangle$ which is a 
2-generator 1-relator group with non-trivial center since $a^4=u^9$,
thus recovering the form of $G$ in Case (ii).

We finish with a question:
\begin{question}
  With the exception of cyclic groups and $BS(1,n)$, does every 1-relator
  group or even group of deficiency 1 have infinite dimensional second bounded
  cohomology?
\end{question}
This reduces to the case of strictly ascending HNN extensions of
finitely generated free groups, as well as needing to eliminate the
existence of strictly ascending HNN extensions of other finitely
generated groups for the deficiency one version of this question.

\begin{thebibliography}{99}

\bibitem{bmrn} G.\,Baumslag,
  Residual nilpotence and relations in free groups,
  {\it J. Algebra} {\bf 2} (1965) 271--282.
  
\bibitem{bp} B.\,Baumslag and S.\,J.\,Pride,
Groups with two more generators than relators,
{\it J. London Math. Soc.}. {\bf 17} (1978) 425--426.


\bibitem{bsh} G.\, Baumslag and P.\,B\,Shalen,
Amalgamated products and finitely presented groups,
{\it Comment. Math. Helv.} {\bf 65} (1990) 243--254. 
  
\bibitem{bstfgn} M.\,Bestvina and K.\,Fujiwara,
  Bounded cohomology of subgroups of mapping class groups,
  {\it Geom. Topol.} {\bf 6} (2002) 69--89. 
  
\bibitem{bi} R.\,Bieri,
Deficiency and the geometric invariants of a group.
{\it J. Pure Appl. Algebra}. {\bf 208} (2007), 951--959.

\bibitem{bistr} R.\,Bieri and R.\,Strebel,
Almost finitely presented soluble groups.
{\it Comment. Math. Helv.}. {\bf 53} (1978), 258--278.

\bibitem{bs} A.\,Borisov and M.\,Sapir,
Polynomial maps over finite fields and residual finiteness of mapping
tori of group endomorphisms,
{\it Invent. Math.} {\bf 160} (2005) 341--356.

\bibitem{brsp2}
  A.\,Borisov and M.\,Sapir,
  Polynomial maps over p-adics and residual properties of mapping tori
  of group endomorphisms, {\it Int. Math. Res. Not.}
  IMRN 2009, no. 16, 3002--3015.

\bibitem{mearx} J.\,O.\,Button,
{\it Large mapping tori of free group endomorphisms},
\texttt{http://arxiv.org/0511715}

  
\bibitem{med1} J.\,O.\,Button,
Large groups of deficiency 1.
{\it Israel J. Math.}. {\bf 167} (2008), 111--140.

\bibitem{me1r} J.\,O.\,Button,
Largeness of LERF and 1-relator groups,
{\it Groups Geom. Dyn.} {\bf 4} (2010), 709--738.

\bibitem{caldun} D.\,Calegari and N.\,M.\,Dunfield,
  An ascending HNN extension of a free group inside $SL_2(\Cc)$,
  {\it Proc. Amer. Math. Soc.} {\bf 134} (2006) 3131--3136. 

\bibitem{invt} P.\,Corvaja, A.\,S.\,Rapinchuk, J.\,Ren, U.\,M.\,Zannier,
  Non-virtually abelian anisotropic linear groups are not boundedly generated.
  {\it Invent. Math.} {\bf 227} (2022) 1--26. 

\bibitem{dm?} F.\,Dahmani and R.\,Li,
  Relative hyperbolicity for automorphisms of free products and free groups,
  {\it J. Topol. Anal.} {\bf 14} (2022) 55--92. 
  
\bibitem{dnwdybk}
  W.\,Dicks and M.\,J.\,Dunwoody,
  {\it Groups acting on graphs}. Cambridge Studies in Advanced Mathematics, 17.
  Cambridge University Press, Cambridge, 1989.


\bibitem{drusap} C.\,Dru\c tu and M.\,Sapir,
Non-linear residually finite groups,
{\it J. Algebra} {\bf 284} (2005).
  
\bibitem{epfuj} D.\,B.\,A.\,Epstein and K.\,Fujiwara,
The second bounded cohomology of word-hyperbolic groups,
  {\it Topology} {\bf 36} (1997) 1275--1289.


\bibitem{flm} B.\,Farb, A.\,Lubotzky and Y.\,Minsky,
Rank-1 phenomena for mapping class groups,
{\it Duke Math. J.} {\bf 106} (2001) 581–-597.

\bibitem{frdfrm} M.\,Forester,
Deformation and rigidity of simplicial group actions on trees,
 {\it Geom. Topol.} {\bf 6} (2002) 219--267.

 
\bibitem{kj2} K.\,Fujiwara,
The second bounded cohomology of an amalgamated free product of groups,
{\it Trans. Amer. Math. Soc.} {\bf 352} (2000) 1113--1129. 

\bibitem{ghst} P.\,Ghosh,
  Relative hyperbolicity of free-by-cyclic extensions,
  {\it Compos. Math.} {\bf 159} (2023) 153--183.


\bibitem{grig} R.\,I.\,Grigorchuk,
Bounded cohomology of group constructions. (Russian)
{\it Mat. Zametki} {\bf 59} (1996) 546--550; translation in
{\it Math. Notes} {\bf 59} (1996) 392--394. 


\bibitem{hgws} M.\,F.\,Hagen and D.\,T.\,Wise,
Cubulating hyperbolic free-by-cyclic groups:
the irreducible case,
{\it Duke Math. J.} {\bf 165} (2016) 1753--1813.


\bibitem{hll} J.\,A.\,Hillman,
  Deficiencies of lattices in connected Lie groups,
  {\it Bull. Austral. Math. Soc.} {\bf 65} (2002) 393--397.


\bibitem{hffmstrs} F.\,Hofmann,
  Quasimorphisms and bounded cohomology,
  Master thesis, Universit\"at Regensburg (2023).
  
\bibitem{sstocan} A.\,Iozzi, C.\,Pagliantini, A.\,Sisto,
  Characterising actions on trees yielding non-trivial quasimorphisms,
  {\it Ann. Math. Qué.} {\bf 45} (2021) 185--202.
  
\bibitem{dss} D.\,H.\,Kochloukova,
Some Novikov rings that are von Neumann finite and knot-like groups.
{\it Comment. Math. Helv.}. {\bf 81} (2006), 931--943.

\bibitem{krp} P.\,H.\,Kropholler,
Baumslag-Solitar groups and some other groups of
cohomological dimension two,
{\it Comment. Math. Helvetici} {\bf 65} (1990) 547--558.

\bibitem{lac} M.\,Lackenby,
  New lower bounds on subgroup growth and homology growth,
  {\it Proc. Lond. Math. Soc.} {\bf 98} (2009) 271--297.
  
\bibitem{lvrf} G.\,Levitt,
Generalized Baumslag-Solitar groups: rank and finite index subgroups,
{\it Ann. Inst. Fourier (Grenoble)} {\bf 65} (2015) 725--762.

\bibitem{lgv3}
  A.\,Logan,
{\it The conjugacy problem for ascending HNN-extensions of free groups},
\texttt{http://arxiv.org/2209.04357v3}

\bibitem{lgv4}
  A.\,Logan,
The conjugacy problem for ascending HNN-extensions of free groups,
{\it to appear}.

\bibitem{ltt} J.\,Lott,
  Deficiencies of lattice subgroups of Lie groups,
  {\it Bull. London Math. Soc.} {\bf 31} (1999) 191--196.


  
\bibitem{lbsgcg} A.\,Lubotzky,
  Alexander Group presentation, $p$-adic analytic groups and lattices in
  $SL_2(\Cc)$,
  {\it Ann. of Math.} {\bf 118} (1983) 115--130. 

\bibitem{lssg} A.\,Lubotzky and D.\,Segal, {\it Subgroup growth}.
Progress in Mathematics 212, Birkha\"user Verlag, Basel, 2003.



\bibitem{ls} R.\,C.\,Lyndon and P.\,E.\,Schupp,
{\it Combinatorial Group Theory}. Springer-Verlag, Berlin-Heidelberg-New
York 1977.



\bibitem{motr} A.\,Minaysan and D.\,Osin,
  Acylindrical hyperbolicity of groups acting on trees,
  {\it Math. Ann.} {\bf 362} (2015) 1055--1105. 


  
\bibitem{mofup} A.\,Minaysan and D.\,Osin,
  Correction to: Acylindrical hyperbolicity of groups acting on trees,
  {\it Math. Ann.} {\bf 373} (2019) 895--900.
  

\bibitem{mur} A.\,Muranov, 
On torsion-free groups with finite regular file bases,
{\it Trans. Amer. Math. Soc.} {\bf 359} (2007) 3609--3645. 


\bibitem{mut} J.\,P.\,Mutanguha,
  The dynamics and geometry of free group endomorphisms,
  {\it Adv. Math.} 384 (2021), Paper No. 107714, 60 pp.

\bibitem{pyb} L.\,Pyber,
  Bounded generation and subgroup growth,
  {\it Bull. London Math. Soc.} {\bf  34} (2002), 55--60. 

\bibitem{pysg} L.\,Pyber and D.\,Segal,
Finitely generated groups with polynomial index growth,
{\it J. Reine Angew. Math.} {\bf 612} (2007), 173--211.


\bibitem{gz2} G.\,Wilkes,
  {\it   Profinite groups and residual finiteness},
  EMS Textbooks in Mathematics. EMS Press, Berlin, 2024.
\end{thebibliography}

\end{document}
